%% 第一章

\chapter{绪论}
\chaptermark{绪论}

\section{研究背景与意义}

近些年来，多模态大模型在视觉问答、图像理解、视频理解以及复杂推理等方面都取得了快速发展。随着模型视觉感知、语言理解和指令跟随能力的增强，研究者开始进一步关注其在具身智能体中的应用潜力，例如机器人操作、导航、无人机巡检和低空目标搜索等任务\cite{shao2026uavwithllm,zeng2023robotics,li2025perception,han2025survey,lin2024embodied}。与静态图像或者普通视频理解不同的是，具身智能体的视觉输入是由自身运动产生的，观测内容会随着位置、姿态、视角以及任务行为的变化而不断变化。因此，仅仅用通用图像问答指标来衡量模型是否有可靠的具身推理能力是不合适的。

无人机属于典型的主动观测智能体。低空无人机在开放三维空间中飞行时，由于高度、朝向、距离以及遮挡状况的变化，会致使观测对象出现显著外观差异\cite{shao2025uavvln}。对无人机智能体而言，完成可靠推理不只是识别出外部的地标或者描述出场景，还要理解自身相对于地标的位置，未来的视角变化，当前所处的飞行动作以及时间过程。无人机推理要依靠对自身状态以及外界环境的把握，这样的能力结构要比传统的被动视觉理解更为复杂。

现有无人机场景数据集和评测任务已经覆盖空中视觉语言导航、城市空间问答和低空目标搜索等方向\cite{liu2023aerialvln,lee2025citynav,zhao2025cityeqa,ji2026towards}。与此同时，空中视频理解、事件识别和目标检测类数据集也为无人机场景建模提供了重要基础\cite{li2021uavhuman,zhao2025urbanvideobench,barekatain2017okutama,du2018unmanned}。这些研究为无人机多模态任务奠定了基础，但多数工作仍以环境理解、导航成功率或事件识别为主要目标，对无人机自身状态的显式评测不足。与此同时，许多多模态大模型综合评测虽然能够衡量视觉问答、跨学科理解和空间关系推理能力\cite{fu2023mme,liu2024mmbench,yu2023mmvet,yue2024mmmu,ge2025mllm}，但其输入通常来自静态图像或弱时序视频，与无人机主动运动产生的多视角、连续时空观测仍存在差距\cite{xu2026citycube}。

就能力优化而言，多模态大模型向具身场景推进的时候，研究重心由原来的扩大模型规模转变为面向任务目标的后训练优化。监督微调（Supervised Fine-Tuning, SFT）可以利用高质量样本使模型学会稳定的输入输出映射，而低秩适配（Low-Rank Adaptation, LoRA）、QLoRA 等参数高效微调方法又大大降低了多模态模型在有限算力下的适配成本。另外，基于可验证奖励的强化后训练方法被用来处理结构化推理、数学推理以及工具调用等任务，从而使模型可以针对格式合法性、答案正确性或者证据质量进行定向校正。无人机双认知推理能力优化不能只追求自然语言回答更加流畅，而应该从答案是否正确、空间或者时间证据是否可靠、不同认知分支是否协同等几个方面展开。因此，本文在完成能力测评之后，进一步创建图像任务训练集，并将 SFT 与联合群组相对策略优化（Group Relative Policy Optimization, GRPO）结合，检验证据感知训练对于无人机多模态双认知能力的加强效果。

基于上述背景，本文围绕“多模态大模型是否能够在无人机主动观测场景中同时理解自身和外界，并给出可靠的空间或时间证据”这一问题展开研究。本文设计并构建的 UAV-DualCog 测试基准将无人机具身推理能力划分为自我状态认知和环境情况认知两条主线，在图像和视频两种媒介下设计六类任务，并将目标框和时间区间作为正式评测对象。该设计能够区分“语义判断正确但证据不足”和“语义与证据同时正确”两类情况，从而更细致地刻画多模态大模型在无人机场景中的真实能力。通过大规模、系统性评测确定当前多模态大模型“自我状态认知不足、语义与证据脱节和跨媒介稳定性”的问题后，本文构建 UAV-DualCog-Train 图像任务训练集，以 Qwen 3.5 4B 为基座模型进行包含 SFT 和联合 GRPO 的两阶段微调，显著提升了模型的图像任务能力，并在视频任务中验证了部分跨媒介迁移效果。

\section{研究目标与主要内容}

本文的总体目标是完成一套面向无人机双认知时空推理的多模态大模型能力评测与优化方法研究。研究线路分为两个相互衔接且并列展开的部分：第一，构建 UAV-DualCog 测试基准并系统评估当前多模态大模型在无人机双认知任务中的能力结构；第二，根据评测中暴露出的语义与证据脱节、自我状态认知不足和跨媒介稳定性不足等问题，构建 UAV-DualCog-Train 训练集，并实现基于监督微调和 GRPO 的能力优化方法。论文正文完整呈现能力测试方法与实验、能力优化方法与实验，以及两部分之间的证据关联。

围绕上述目标，本文完成的工作可以概括为三个方面。

第一，构建面向无人机具身推理的双认知能力体系。本文将无人机多模态推理拆分为自我状态认知和环境情况认知两类顶层能力：前者关注无人机对自身位置、未来视角、飞行行为和行为时间区间的理解，后者关注无人机对目标方位、动作决策、地标可见次数和可见区间的理解。该体系为后续任务设计、数据构建和实验分析提供统一能力框架。

第二，完成 UAV-DualCog 评测任务、测试集构建、系统评测及结果分析。本文设计四类图像任务和两类视频任务，形成“能力维度、观测媒介、证据形式”相结合的三维评测结构；同时实现四阶段自动化数据构建工具链，以 AerialVLN 仿真场景为基础完成场景扫描与点云融合、地标资产构建与语义标注、行为驱动视频任务生成和证据感知图像任务生成。在此基础上，本文对闭源模型、开源模型和空间推理微调模型进行统一评测，分析语义正确率、空间定位质量、时间定位质量、跨媒介一致性和双认知平衡性等能力结构。

第三，完成训练集构建、能力优化方法设计、优化实验及结果分析。根据评测中暴露出的自我状态认知缺乏、语义与证据脱离、跨媒介稳定性欠缺等状况，本文构建与测试场景隔离的图像任务训练集，并以 Qwen 3.5 4B 为基座模型，经由多模态监督微调并围绕格式、语义、空间证据展开联合 GRPO 来进行能力优化。训练完成后，本文进一步对优化模型进行图像任务直接收益和视频任务迁移效果分析。

\section{本文主要贡献}

本文的主要贡献可以概括为以下三点。

首先，本文面向无人机场景设计并构建了双认知时空推理评测框架。不同于现有无人机视觉语言导航、目标搜索或通用视觉问答评测主要关注外部环境理解和任务完成结果，该框架将“理解自身”和“理解外界”同时纳入顶层能力设计，使评测更加贴近具身无人机的真实推理需求。

其次，本文实现了从场景几何数据到图像、视频问答样本的完整数据集构建链路。与只依赖人工收集图像或视频的静态数据集不同，该链路以 AerialVLN 仿真场景、语义点云、地标复核、层次化飞行行为和结构化任务导出为基础，形成包含 12 个测试场景、512 个有效地标、4096 个图像任务样本和 2048 个视频任务样本的稳定测试集。

最后，本文形成了从能力评测到能力优化的研究闭环。实验表明，当前模型在无人机场景中能够捕捉部分目标、方位和动态线索，但是还没有形成稳定的、均衡的、可以被证据验证的双认知能力。尤其是自我状态认知明显低于环境情况认知，语义正确性也不能直接转化为空间或者时间的证据可靠性。基于这些发现，本文进一步实现了包含 SFT 和联合 GRPO 的两阶段证据感知训练路线，并对优化模型进行任务级评测。结果显示，优化模型在图像任务上显著提升答案正确性和空间证据质量，在视频任务上的全局准确率也有所提升。


\section{论文组织结构}

本文共分为五章。

第一章为绪论，介绍研究背景、研究意义、研究目标、主要内容和论文结构。

第二章介绍相关研究与理论基础，重点梳理无人机具身智能与低空多模态评测、多模态大模型综合评测与空间推理、能力优化训练方法等方面的工作。

第三章为 UAV-DualCog 双认知能力测试方法与实验，介绍测试基准设计、数据构建工具链、模型评测设置、图像任务结果、视频任务结果、跨媒介迁移分析和双认知差异分析。

第四章为能力优化方法与实验，介绍 UAV-DualCog-Train 构建方案、多模态监督微调方法、联合 GRPO 优化方法、训练过程以及优化模型评测结果。

第五章为总结与展望，总结本文已完成工作、主要实验发现、研究局限以及后续可进一步完善的方向。
